\section{Proposed Postdoctoral Project}

\begin{frame}
\frametitle{Aims}
\framesubtitle{An autonomous discovery loop for halide solid electrolytes}

\noindent
\begin{minipage}[t]{0.315\linewidth}
\begin{tcolorbox}[colback=primary!6,colframe=primary,title={\small Aim 1},fonttitle=\bfseries,coltitle=white,height=4.7cm]
\small Extend IonForge to halide chemistries (Cl, Br, I frameworks) beyond the oxide and sulfide systems covered in my thesis.
\end{tcolorbox}
\end{minipage}\hfill
\begin{minipage}[t]{0.315\linewidth}
\begin{tcolorbox}[colback=secondary!10,colframe=secondary,title={\small Aim 2},fonttitle=\bfseries,coltitle=white,height=4.7cm]
\small Couple IonForge screening with the Nimbus AutoSynth robotic platform for closed-loop candidate validation.
\end{tcolorbox}
\end{minipage}\hfill
\begin{minipage}[t]{0.315\linewidth}
\begin{tcolorbox}[colback=accent!12,colframe=accent,title={\small Aim 3},fonttitle=\bfseries,coltitle=black,height=4.7cm]
\small Release an open benchmark dataset for halide solid electrolytes for the wider community.
\end{tcolorbox}
\end{minipage}

\vspace{0.5cm}
\begin{itemize}
  \item Target output: 5--8 experimentally validated halide conductor candidates within the fellowship period
  \item Deliverable: IonForge-Halide, released alongside the benchmark dataset
\end{itemize}
\end{frame}

% -----------------------------------------------------------------------------
\begin{frame}
\frametitle{24-Month Timeline}
\framesubtitle{Proposed schedule for the fellowship period}

\begin{center}
\resizebox{0.98\linewidth}{!}{%
\begin{tikzpicture}[x=1.0cm,y=1cm,font=\scriptsize]
\foreach \q in {0,...,8}{
  \draw[gray!35] (\q,0) -- (\q,-4.3);
}
\foreach \q/\lab in {0/Q1,1/Q2,2/Q3,3/Q4,4/Q5,5/Q6,6/Q7,7/Q8}{
  \node[above] at (\q+0.5,0) {\lab};
}
\node[anchor=east,text width=3.4cm,align=right] at (-0.15,-0.5) {Extend IonForge to halides};
\draw[fill=primary,draw=primary] (0,-0.9) rectangle (3,-0.1);
\node[anchor=east,text width=3.4cm,align=right] at (-0.15,-1.5) {Integrate with Nimbus AutoSynth};
\draw[fill=secondary,draw=primary] (1,-1.9) rectangle (5,-1.1);
\node[anchor=east,text width=3.4cm,align=right] at (-0.15,-2.5) {Closed-loop discovery campaigns};
\draw[fill=accent,draw=primary] (3,-2.9) rectangle (7,-2.1);
\node[anchor=east,text width=3.4cm,align=right] at (-0.15,-3.5) {Dataset release \& manuscripts};
\draw[fill=highlight,draw=primary] (5,-3.9) rectangle (8,-3.1);
\end{tikzpicture}}
\end{center}

\vspace{0.3cm}
\small\color{secondary} Bars indicate active phases; overlaps reflect iterative, not sequential, execution.
\end{frame}

% =============================================================================
